\documentclass[aspectratio=169]{beamer}
\usetheme{Madrid}
\usecolortheme{default}
\usepackage[utf8]{inputenc}
\usepackage[spanish,es-minimal]{babel}
\usepackage{amsmath}
\usepackage{amssymb}
\usepackage{booktabs}
\usepackage{multirow}
\usepackage{graphicx}

\definecolor{unsblue}{RGB}{30,60,120}
\setbeamercolor{structure}{fg=unsblue}
\setbeamercolor{palette primary}{bg=unsblue,fg=white}
\setbeamercolor{palette secondary}{bg=unsblue!80,fg=white}
\setbeamercolor{palette tertiary}{bg=unsblue!60,fg=white}

\setbeamertemplate{navigation symbols}{}
\setbeamertemplate{footline}[frame number]

\title[MAD1 - Clase 22]{Métodos de Análisis de Datos 1}
\subtitle{Clase 22: ANOVA de Dos Factores y Diseños Experimentales}
\author{Prof. José}
\institute{Departamento de Matemática \\ Universidad Nacional del Sur}
\date{}

\begin{document}

\begin{frame}
\titlepage
\end{frame}

\begin{frame}{Contenidos de la Clase}
\tableofcontents
\end{frame}

\section{ANOVA de Dos Factores}

\begin{frame}{Motivación}
\begin{itemize}
\item Hasta ahora: un solo factor de variación (ANOVA de un factor)
\item En la práctica: múltiples factores pueden afectar la respuesta simultáneamente
\item Ejemplos:
\begin{itemize}
\item Rendimiento de cultivos: tipo de fertilizante Y cantidad de agua
\item Resistencia de materiales: temperatura Y presión
\item Efectividad de tratamiento: dosis Y edad del paciente
\end{itemize}
\item ANOVA de dos factores permite:
\begin{itemize}
\item Estudiar efecto de cada factor por separado
\item Detectar interacciones entre factores
\item Mayor eficiencia experimental
\end{itemize}
\end{itemize}
\end{frame}

\begin{frame}{Modelo de ANOVA de Dos Factores}
\begin{block}{Modelo con interacción}
$$Y_{ijk} = \mu + \alpha_i + \beta_j + (\alpha\beta)_{ij} + \varepsilon_{ijk}$$
\end{block}

Donde:
\begin{itemize}
\item $Y_{ijk}$: observación $k$ del factor A nivel $i$, factor B nivel $j$
\item $\mu$: media global
\item $\alpha_i$: efecto del factor A en el nivel $i$ ($i=1,\ldots,a$)
\item $\beta_j$: efecto del factor B en el nivel $j$ ($j=1,\ldots,b$)
\item $(\alpha\beta)_{ij}$: efecto de interacción entre A y B
\item $\varepsilon_{ijk}$: error aleatorio, $\varepsilon_{ijk} \sim N(0,\sigma^2)$
\item $k=1,\ldots,n$ réplicas en cada combinación $(i,j)$
\end{itemize}
\end{frame}

\begin{frame}{Restricciones del Modelo}
Para que el modelo sea identificable, se imponen restricciones:

$$\sum_{i=1}^{a} \alpha_i = 0, \quad \sum_{j=1}^{b} \beta_j = 0$$

$$\sum_{i=1}^{a} (\alpha\beta)_{ij} = 0 \text{ para todo } j, \quad \sum_{j=1}^{b} (\alpha\beta)_{ij} = 0 \text{ para todo } i$$

\vspace{0.3cm}
Esto asegura que los efectos sean desviaciones respecto a la media global.
\end{frame}

\begin{frame}{Concepto de Interacción}
\begin{block}{¿Qué es la interacción?}
La interacción ocurre cuando el efecto de un factor depende del nivel del otro factor.
\end{block}

\begin{columns}
\begin{column}{0.5\textwidth}
\textbf{Sin interacción:}
\begin{itemize}
\item Efectos aditivos
\item Líneas paralelas en gráfico
\item $(\alpha\beta)_{ij} = 0$ para todo $i,j$
\end{itemize}
\end{column}
\begin{column}{0.5\textwidth}
\textbf{Con interacción:}
\begin{itemize}
\item Efectos no aditivos
\item Líneas no paralelas
\item $(\alpha\beta)_{ij} \neq 0$ para algún $i,j$
\end{itemize}
\end{column}
\end{columns}

\vspace{0.3cm}
Si hay interacción significativa, no tiene sentido interpretar los efectos principales por separado.
\end{frame}

\begin{frame}{Tabla ANOVA de Dos Factores}
\small
\begin{table}
\begin{tabular}{lccccc}
\toprule
\textbf{Fuente} & \textbf{SC} & \textbf{gl} & \textbf{CM} & \textbf{F} & \textbf{p-valor} \\
\midrule
Factor A & $SC_A$ & $a-1$ & $CM_A = \frac{SC_A}{a-1}$ & $F_A = \frac{CM_A}{CM_E}$ & $P(F > F_A)$ \\[0.3cm]
Factor B & $SC_B$ & $b-1$ & $CM_B = \frac{SC_B}{b-1}$ & $F_B = \frac{CM_B}{CM_E}$ & $P(F > F_B)$ \\[0.3cm]
Interacción AB & $SC_{AB}$ & $(a-1)(b-1)$ & $CM_{AB} = \frac{SC_{AB}}{(a-1)(b-1)}$ & $F_{AB} = \frac{CM_{AB}}{CM_E}$ & $P(F > F_{AB})$ \\[0.3cm]
Error & $SC_E$ & $ab(n-1)$ & $CM_E = \frac{SC_E}{ab(n-1)}$ & & \\[0.3cm]
\midrule
Total & $SC_T$ & $abn-1$ & & & \\
\bottomrule
\end{tabular}
\end{table}
\end{frame}

\begin{frame}{Cálculo de Sumas de Cuadrados}
\begin{align*}
SC_T &= \sum_{i=1}^{a}\sum_{j=1}^{b}\sum_{k=1}^{n} (Y_{ijk} - \bar{Y}_{\cdot\cdot\cdot})^2 \\[0.3cm]
SC_A &= bn\sum_{i=1}^{a} (\bar{Y}_{i\cdot\cdot} - \bar{Y}_{\cdot\cdot\cdot})^2 \\[0.3cm]
SC_B &= an\sum_{j=1}^{b} (\bar{Y}_{\cdot j\cdot} - \bar{Y}_{\cdot\cdot\cdot})^2 \\[0.3cm]
SC_{AB} &= n\sum_{i=1}^{a}\sum_{j=1}^{b} (\bar{Y}_{ij\cdot} - \bar{Y}_{i\cdot\cdot} - \bar{Y}_{\cdot j\cdot} + \bar{Y}_{\cdot\cdot\cdot})^2 \\[0.3cm]
SC_E &= SC_T - SC_A - SC_B - SC_{AB}
\end{align*}
\end{frame}

\begin{frame}{Hipótesis y Pruebas}
\textbf{Tres pruebas independientes:}

\begin{enumerate}
\item \textbf{Interacción:} $H_0: (\alpha\beta)_{ij} = 0$ para todo $i,j$
\begin{itemize}
\item Estadístico: $F_{AB} = \frac{CM_{AB}}{CM_E} \sim F_{(a-1)(b-1), ab(n-1)}$
\end{itemize}

\vspace{0.2cm}

\item \textbf{Efecto principal de A:} $H_0: \alpha_1 = \alpha_2 = \cdots = \alpha_a = 0$
\begin{itemize}
\item Estadístico: $F_A = \frac{CM_A}{CM_E} \sim F_{a-1, ab(n-1)}$
\end{itemize}

\vspace{0.2cm}

\item \textbf{Efecto principal de B:} $H_0: \beta_1 = \beta_2 = \cdots = \beta_b = 0$
\begin{itemize}
\item Estadístico: $F_B = \frac{CM_B}{CM_E} \sim F_{b-1, ab(n-1)}$
\end{itemize}
\end{enumerate}

\vspace{0.2cm}
\textbf{Importante:} Si la interacción es significativa, interpretar con cuidado los efectos principales.
\end{frame}

\begin{frame}{Ejemplo: Rendimiento de Cultivos}
\small
Un agrónomo estudia el rendimiento (kg/ha) según dos factores:
\begin{itemize}
\item Factor A (Fertilizante): 3 tipos (F1, F2, F3)
\item Factor B (Riego): 2 niveles (Bajo, Alto)
\item 4 réplicas por combinación (total: 24 observaciones)
\end{itemize}

\begin{table}
\centering
\begin{tabular}{lccc}
\toprule
\textbf{Riego} & \textbf{F1} & \textbf{F2} & \textbf{F3} \\
\midrule
Bajo & 45, 47, 46, 48 & 52, 54, 53, 55 & 48, 50, 49, 51 \\
Alto & 55, 57, 56, 58 & 58, 60, 59, 61 & 62, 64, 63, 65 \\
\bottomrule
\end{tabular}
\end{table}

¿Hay efecto del fertilizante? ¿Del riego? ¿Interacción entre ambos?
\end{frame}

\begin{frame}{Ejemplo: Medias Marginales}
\begin{table}
\centering
\begin{tabular}{lcccc}
\toprule
\textbf{Riego} & \textbf{F1} & \textbf{F2} & \textbf{F3} & \textbf{Media} \\
\midrule
Bajo & 46.5 & 53.5 & 49.5 & 49.83 \\
Alto & 56.5 & 59.5 & 63.5 & 59.83 \\
\midrule
\textbf{Media} & 51.5 & 56.5 & 56.5 & 54.83 \\
\bottomrule
\end{tabular}
\end{table}

\vspace{0.3cm}
Observaciones preliminares:
\begin{itemize}
\item Riego alto produce mayor rendimiento (59.83 vs 49.83)
\item F2 y F3 parecen mejores que F1
\item ¿Hay interacción? Necesitamos el gráfico y la prueba formal
\end{itemize}
\end{frame}

\section{Diseños Experimentales Clásicos}

\begin{frame}{¿Por qué Diseñar Experimentos?}
\begin{block}{Diseño experimental}
Conjunto de técnicas para planificar experimentos de manera que los datos obtenidos permitan conclusiones válidas y objetivas.
\end{block}

\vspace{0.3cm}
\textbf{Objetivos:}
\begin{itemize}
\item Maximizar información obtenida con recursos limitados
\item Controlar variabilidad no deseada
\item Garantizar validez de las conclusiones
\item Permitir estimación del error experimental
\end{itemize}

\vspace{0.3cm}
\textbf{Principios básicos:}
\begin{enumerate}
\item Réplica: repetir observaciones para estimar error
\item Aleatorización: asignar tratamientos al azar
\item Control local: bloquear fuentes de variación conocidas
\end{enumerate}
\end{frame}

\begin{frame}{Clasificación de Diseños}
\begin{enumerate}
\item \textbf{Diseño Completamente Aleatorizado (DCA):}
\begin{itemize}
\item Asignación completamente al azar
\item Unidades experimentales homogéneas
\item Más simple, menos eficiente si hay heterogeneidad
\end{itemize}

\vspace{0.2cm}

\item \textbf{Diseño en Bloques Completos Aleatorizados (DBCA):}
\begin{itemize}
\item Agrupa unidades en bloques homogéneos
\item Aleatoria dentro de cada bloque
\item Controla variación entre bloques
\end{itemize}

\vspace{0.2cm}

\item \textbf{Diseño en Cuadrado Latino (DCL):}
\begin{itemize}
\item Controla dos fuentes de variación simultáneamente
\item Estructura de filas y columnas
\item Muy eficiente pero restrictivo
\end{itemize}
\end{enumerate}
\end{frame}

\subsection{Diseño Completamente Aleatorizado (DCA)}

\begin{frame}{DCA: Características}
\textbf{Estructura:}
\begin{itemize}
\item $t$ tratamientos
\item $n_i$ réplicas del tratamiento $i$ ($i=1,\ldots,t$)
\item Total: $N = \sum_{i=1}^{t} n_i$ unidades experimentales
\item Asignación completamente al azar
\end{itemize}

\vspace{0.3cm}
\textbf{Modelo:}
$$Y_{ij} = \mu + \tau_i + \varepsilon_{ij}$$

Donde:
\begin{itemize}
\item $Y_{ij}$: observación $j$ del tratamiento $i$
\item $\tau_i$: efecto del tratamiento $i$
\item $\varepsilon_{ij} \sim N(0,\sigma^2)$ independientes
\end{itemize}
\end{frame}

\begin{frame}{DCA: Tabla ANOVA}
\begin{table}
\centering
\begin{tabular}{lccccc}
\toprule
\textbf{Fuente} & \textbf{SC} & \textbf{gl} & \textbf{CM} & \textbf{F} & \textbf{p-valor} \\
\midrule
Tratamientos & $SC_{Trat}$ & $t-1$ & $CM_{Trat}$ & $F = \frac{CM_{Trat}}{CM_E}$ & $P(F > F_{obs})$ \\[0.2cm]
Error & $SC_E$ & $N-t$ & $CM_E$ & & \\[0.2cm]
\midrule
Total & $SC_T$ & $N-1$ & & & \\
\bottomrule
\end{tabular}
\end{table}

\vspace{0.3cm}
Donde:
\begin{align*}
SC_{Trat} &= \sum_{i=1}^{t} n_i(\bar{Y}_{i\cdot} - \bar{Y}_{\cdot\cdot})^2 \\
SC_E &= \sum_{i=1}^{t}\sum_{j=1}^{n_i} (Y_{ij} - \bar{Y}_{i\cdot})^2 \\
SC_T &= \sum_{i=1}^{t}\sum_{j=1}^{n_i} (Y_{ij} - \bar{Y}_{\cdot\cdot})^2
\end{align*}
\end{frame}

\begin{frame}{DCA: Ventajas y Desventajas}
\begin{columns}
\begin{column}{0.5\textwidth}
\textbf{Ventajas:}
\begin{itemize}
\item Muy flexible
\item Permite diferente número de réplicas por tratamiento
\item Análisis simple
\item Máximos grados de libertad para el error
\item Fácil manejo de datos faltantes
\end{itemize}
\end{column}
\begin{column}{0.5\textwidth}
\textbf{Desventajas:}
\begin{itemize}
\item Requiere unidades experimentales homogéneas
\item Ineficiente si hay fuentes de variación conocidas
\item Mayor variabilidad del error
\end{itemize}
\end{column}
\end{columns}

\vspace{0.4cm}
\textbf{Cuándo usar:}
\begin{itemize}
\item Material experimental homogéneo
\item Pocas unidades experimentales
\item Condiciones controladas (laboratorio)
\end{itemize}
\end{frame}

\subsection{Diseño en Bloques Completos Aleatorizados (DBCA)}

\begin{frame}{DBCA: Motivación}
\textbf{Problema del DCA:}
\begin{itemize}
\item Si las unidades experimentales son heterogéneas, aumenta la variabilidad del error
\item Reduce potencia de las pruebas
\end{itemize}

\vspace{0.3cm}
\textbf{Solución del DBCA:}
\begin{itemize}
\item Agrupar unidades similares en bloques
\item Cada bloque contiene todos los tratamientos (una vez cada uno)
\item Aleatorización dentro de cada bloque
\item Eliminar variabilidad entre bloques del error
\end{itemize}

\vspace{0.3cm}
\textbf{Ejemplos de bloques:}
\begin{itemize}
\item Lotes de terreno con características similares
\item Días de experimentación
\item Lotes de materia prima
\item Pares de sujetos emparejados
\end{itemize}
\end{frame}

\begin{frame}{DBCA: Estructura}
\textbf{Componentes:}
\begin{itemize}
\item $t$ tratamientos
\item $b$ bloques
\item Una observación por combinación tratamiento-bloque
\item Total: $N = tb$ observaciones
\end{itemize}

\vspace{0.3cm}
\textbf{Modelo:}
$$Y_{ij} = \mu + \tau_i + \beta_j + \varepsilon_{ij}$$

Donde:
\begin{itemize}
\item $Y_{ij}$: observación del tratamiento $i$ en el bloque $j$
\item $\mu$: media global
\item $\tau_i$: efecto del tratamiento $i$ ($i=1,\ldots,t$)
\item $\beta_j$: efecto del bloque $j$ ($j=1,\ldots,b$)
\item $\varepsilon_{ij} \sim N(0,\sigma^2)$ independientes
\end{itemize}
\end{frame}

\begin{frame}{DBCA: Tabla ANOVA}
\begin{table}
\centering
\begin{tabular}{lccccc}
\toprule
\textbf{Fuente} & \textbf{SC} & \textbf{gl} & \textbf{CM} & \textbf{F} & \textbf{p-valor} \\
\midrule
Tratamientos & $SC_{Trat}$ & $t-1$ & $CM_{Trat}$ & $F_{Trat} = \frac{CM_{Trat}}{CM_E}$ & $P(F > F_{Trat})$ \\[0.2cm]
Bloques & $SC_{Bloq}$ & $b-1$ & $CM_{Bloq}$ & $F_{Bloq} = \frac{CM_{Bloq}}{CM_E}$ & $P(F > F_{Bloq})$ \\[0.2cm]
Error & $SC_E$ & $(t-1)(b-1)$ & $CM_E$ & & \\[0.2cm]
\midrule
Total & $SC_T$ & $tb-1$ & & & \\
\bottomrule
\end{tabular}
\end{table}

\vspace{0.3cm}
Donde:
\begin{align*}
SC_{Trat} &= b\sum_{i=1}^{t} (\bar{Y}_{i\cdot} - \bar{Y}_{\cdot\cdot})^2, \quad
SC_{Bloq} = t\sum_{j=1}^{b} (\bar{Y}_{\cdot j} - \bar{Y}_{\cdot\cdot})^2 \\
SC_E &= SC_T - SC_{Trat} - SC_{Bloq}
\end{align*}
\end{frame}

\begin{frame}{DBCA: Hipótesis}
\textbf{Prueba principal (tratamientos):}
$$H_0: \tau_1 = \tau_2 = \cdots = \tau_t = 0$$
$$H_1: \text{Al menos un } \tau_i \neq 0$$

Estadístico: $F_{Trat} = \frac{CM_{Trat}}{CM_E} \sim F_{t-1, (t-1)(b-1)}$ bajo $H_0$

\vspace{0.3cm}
\textbf{Prueba secundaria (bloques):}
$$H_0: \beta_1 = \beta_2 = \cdots = \beta_b = 0$$

Estadístico: $F_{Bloq} = \frac{CM_{Bloq}}{CM_E} \sim F_{b-1, (t-1)(b-1)}$ bajo $H_0$

\vspace{0.3cm}
\textbf{Nota:} La prueba de bloques indica si el bloqueo fue efectivo, pero no es el objetivo principal del análisis.
\end{frame}

\begin{frame}{DBCA: Ventajas y Desventajas}
\begin{columns}
\begin{column}{0.5\textwidth}
\textbf{Ventajas:}
\begin{itemize}
\item Controla variación conocida
\item Mayor precisión (menor $CM_E$)
\item Mayor potencia estadística
\item Diseño muy utilizado en la práctica
\end{itemize}
\end{column}
\begin{column}{0.5\textwidth}
\textbf{Desventajas:}
\begin{itemize}
\item Requiere $b \geq t$ para tener grados de libertad del error
\item Menos flexible que DCA
\item Supone ausencia de interacción tratamiento-bloque
\item Datos faltantes complican análisis
\end{itemize}
\end{column}
\end{columns}

\vspace{0.4cm}
\textbf{Cuándo usar:}
\begin{itemize}
\item Material experimental heterogéneo pero agrupable
\item Fuentes de variación conocidas y controlables
\item Experimentos de campo, ensayos clínicos
\end{itemize}
\end{frame}

\begin{frame}{Ejemplo DBCA: Rendimiento de Variedades}
Un investigador compara 4 variedades de trigo en 5 lotes de terreno (bloques). Rendimiento en toneladas/ha:

\vspace{0.2cm}
\small
\begin{table}
\centering
\begin{tabular}{lccccc}
\toprule
\textbf{Variedad} & \textbf{Lote 1} & \textbf{Lote 2} & \textbf{Lote 3} & \textbf{Lote 4} & \textbf{Lote 5} \\
\midrule
V1 & 5.2 & 4.8 & 5.5 & 4.9 & 5.1 \\
V2 & 6.1 & 5.7 & 6.4 & 5.8 & 6.0 \\
V3 & 5.8 & 5.4 & 6.1 & 5.5 & 5.7 \\
V4 & 5.5 & 5.1 & 5.8 & 5.2 & 5.4 \\
\bottomrule
\end{tabular}
\end{table}

\vspace{0.3cm}
\normalsize
¿Hay diferencias significativas entre variedades?
\end{frame}

\subsection{Diseño en Cuadrado Latino (DCL)}

\begin{frame}{Cuadrado Latino: Concepto}
\textbf{Objetivo:} Controlar dos fuentes de variación simultáneamente

\vspace{0.3cm}
\textbf{Estructura:}
\begin{itemize}
\item $t$ tratamientos, $t$ filas, $t$ columnas
\item Cada tratamiento aparece exactamente una vez en cada fila
\item Cada tratamiento aparece exactamente una vez en cada columna
\item Total: $t^2$ unidades experimentales
\end{itemize}

\vspace{0.3cm}
\textbf{Ejemplo con 4 tratamientos (A, B, C, D):}
\begin{table}
\centering
\begin{tabular}{c|cccc}
& Col 1 & Col 2 & Col 3 & Col 4 \\
\hline
Fila 1 & A & B & C & D \\
Fila 2 & B & C & D & A \\
Fila 3 & C & D & A & B \\
Fila 4 & D & A & B & C \\
\end{tabular}
\end{table}
\end{frame}

\begin{frame}{DCL: Modelo y Supuestos}
\textbf{Modelo:}
$$Y_{ijk} = \mu + \tau_i + \rho_j + \gamma_k + \varepsilon_{ijk}$$

Donde:
\begin{itemize}
\item $Y_{ijk}$: observación del tratamiento $i$, fila $j$, columna $k$
\item $\mu$: media global
\item $\tau_i$: efecto del tratamiento $i$
\item $\rho_j$: efecto de la fila $j$
\item $\gamma_k$: efecto de la columna $k$
\item $\varepsilon_{ijk} \sim N(0,\sigma^2)$ independientes
\end{itemize}

\vspace{0.3cm}
\textbf{Supuesto crítico:} No hay interacción entre tratamientos, filas y columnas.
\end{frame}

\begin{frame}{DCL: Tabla ANOVA}
\begin{table}
\centering
\small
\begin{tabular}{lccccc}
\toprule
\textbf{Fuente} & \textbf{SC} & \textbf{gl} & \textbf{CM} & \textbf{F} & \textbf{p-valor} \\
\midrule
Tratamientos & $SC_{Trat}$ & $t-1$ & $CM_{Trat}$ & $F_{Trat} = \frac{CM_{Trat}}{CM_E}$ & $P(F > F_{Trat})$ \\[0.2cm]
Filas & $SC_{Filas}$ & $t-1$ & $CM_{Filas}$ & $F_{Filas} = \frac{CM_{Filas}}{CM_E}$ & $P(F > F_{Filas})$ \\[0.2cm]
Columnas & $SC_{Cols}$ & $t-1$ & $CM_{Cols}$ & $F_{Cols} = \frac{CM_{Cols}}{CM_E}$ & $P(F > F_{Cols})$ \\[0.2cm]
Error & $SC_E$ & $(t-1)(t-2)$ & $CM_E$ & & \\[0.2cm]
\midrule
Total & $SC_T$ & $t^2-1$ & & & \\
\bottomrule
\end{tabular}
\end{table}

\vspace{0.3cm}
Donde:
$$SC_E = SC_T - SC_{Trat} - SC_{Filas} - SC_{Cols}$$
\end{frame}

\begin{frame}{DCL: Ventajas y Desventajas}
\begin{columns}
\begin{column}{0.5\textwidth}
\textbf{Ventajas:}
\begin{itemize}
\item Controla dos fuentes de variación
\item Muy eficiente cuando hay dos gradientes
\item Reduce mucho el error experimental
\end{itemize}
\end{column}
\begin{column}{0.5\textwidth}
\textbf{Desventajas:}
\begin{itemize}
\item Muy restrictivo: $t = $ filas $=$ columnas
\item Pocos grados de libertad para error
\item Supone ausencia de interacciones
\item Datos faltantes son problemáticos
\item Difícil de implementar con muchos tratamientos
\end{itemize}
\end{column}
\end{columns}

\vspace{0.4cm}
\textbf{Cuándo usar:}
\begin{itemize}
\item Dos fuentes de variación conocidas y ortogonales
\item Número pequeño de tratamientos (3-10)
\item Ejemplo: ensayos agrícolas con gradiente norte-sur y este-oeste
\end{itemize}
\end{frame}

\begin{frame}{Ejemplo DCL: Tiempo de Secado}
Se estudia el tiempo de secado (minutos) de una pintura según 4 formulaciones. Se controla:
\begin{itemize}
\item Operador (4 operadores)
\item Día de la semana (4 días)
\end{itemize}

\vspace{0.2cm}
\small
\begin{table}
\centering
\begin{tabular}{lcccc}
\toprule
& \textbf{Lunes} & \textbf{Martes} & \textbf{Miércoles} & \textbf{Jueves} \\
\midrule
Op. 1 & A: 73 & B: 74 & C: 71 & D: 75 \\
Op. 2 & B: 73 & C: 75 & D: 72 & A: 74 \\
Op. 3 & C: 75 & D: 73 & A: 76 & B: 74 \\
Op. 4 & D: 72 & A: 75 & B: 73 & C: 71 \\
\bottomrule
\end{tabular}
\end{table}

\vspace{0.3cm}
\normalsize
¿Hay diferencias entre formulaciones? ¿Hay efecto de operador o día?
\end{frame}

\section{Comparación de Diseños}

\begin{frame}{Eficiencia Relativa}
\textbf{¿Cuál diseño es mejor?}

La eficiencia relativa compara la precisión de dos diseños:

$$\text{Eficiencia Relativa} = \frac{CM_E(\text{Diseño 1})}{CM_E(\text{Diseño 2})} \times \frac{gl_E(\text{Diseño 2}) + 1}{gl_E(\text{Diseño 1}) + 1} \times 100\%$$

\vspace{0.3cm}
\textbf{Interpretación:}
\begin{itemize}
\item ER > 100\%: Diseño 2 es más eficiente que Diseño 1
\item ER = 100\%: Ambos diseños igualmente eficientes
\item ER < 100\%: Diseño 1 es más eficiente que Diseño 2
\end{itemize}

\vspace{0.3cm}
Ejemplo: Si DBCA tiene ER = 130\% respecto a DCA, se necesitarían 30\% más observaciones en DCA para lograr la misma precisión.
\end{frame}

\begin{frame}{Guía para Elegir un Diseño}
\begin{enumerate}
\item \textbf{Material experimental homogéneo, pocos tratamientos:}
\begin{itemize}
\item Usar DCA
\end{itemize}

\vspace{0.2cm}

\item \textbf{Material heterogéneo, una fuente de variación identificable:}
\begin{itemize}
\item Usar DBCA
\end{itemize}

\vspace{0.2cm}

\item \textbf{Dos fuentes de variación identificables y ortogonales:}
\begin{itemize}
\item Usar DCL (si $t$ es pequeño, típicamente $t \leq 10$)
\end{itemize}

\vspace{0.2cm}

\item \textbf{Múltiples factores de interés:}
\begin{itemize}
\item Usar diseños factoriales (ANOVA de dos o más factores)
\end{itemize}

\vspace{0.2cm}

\item \textbf{Recursos muy limitados:}
\begin{itemize}
\item Considerar diseños factoriales fraccionados
\end{itemize}
\end{enumerate}
\end{frame}

\section{Verificación de Supuestos}

\begin{frame}{Supuestos del Modelo}
Todos los diseños ANOVA requieren:

\begin{enumerate}
\item \textbf{Normalidad:} Los errores $\varepsilon_{ij}$ siguen distribución normal
\item \textbf{Homocedasticidad:} Varianza constante $\sigma^2$ para todos los grupos
\item \textbf{Independencia:} Las observaciones son independientes entre sí
\end{enumerate}

\vspace{0.3cm}
\textbf{Consecuencias de violaciones:}
\begin{itemize}
\item Normalidad: Afecta poco con muestras grandes (TLC)
\item Homocedasticidad: Puede invalidar pruebas F
\item Independencia: Problema serio, invalida todo el análisis
\end{itemize}
\end{frame}

\begin{frame}{Verificación de Normalidad}
\textbf{Métodos gráficos:}
\begin{itemize}
\item Gráfico Q-Q de residuos
\item Histograma de residuos
\end{itemize}

\vspace{0.3cm}
\textbf{Pruebas formales:}
\begin{itemize}
\item Shapiro-Wilk (recomendada para $n < 50$)
\item Kolmogorov-Smirnov
\item Anderson-Darling
\end{itemize}

\vspace{0.3cm}
\textbf{Soluciones si se rechaza normalidad:}
\begin{itemize}
\item Transformación de datos (logaritmo, raíz cuadrada, Box-Cox)
\item Pruebas no paramétricas (Kruskal-Wallis, Friedman)
\end{itemize}
\end{frame}

\begin{frame}{Verificación de Homocedasticidad}
\textbf{Métodos gráficos:}
\begin{itemize}
\item Gráfico de residuos vs valores ajustados
\item Gráfico de dispersión por grupo
\end{itemize}

\vspace{0.3cm}
\textbf{Pruebas formales:}
\begin{itemize}
\item Levene (robusta ante desviaciones de normalidad)
\item Bartlett (más potente si hay normalidad)
\item Fligner-Killeen (alternativa robusta)
\end{itemize}

\vspace{0.3cm}
\textbf{Regla práctica:} Si $\frac{s_{max}^2}{s_{min}^2} < 3$, la homocedasticidad es razonable.

\vspace{0.3cm}
\textbf{Soluciones:}
\begin{itemize}
\item Transformación de datos
\item ANOVA de Welch (permite varianzas desiguales)
\end{itemize}
\end{frame}

\begin{frame}{Verificación de Independencia}
\textbf{Causas comunes de dependencia:}
\begin{itemize}
\item Medidas repetidas en el mismo sujeto
\item Observaciones en el tiempo (autocorrelación)
\item Agrupamiento espacial
\item Estructura jerárquica de datos
\end{itemize}

\vspace{0.3cm}
\textbf{Detección:}
\begin{itemize}
\item Gráfico de residuos vs orden de recolección
\item Prueba de Durbin-Watson (autocorrelación)
\item Análisis del diseño experimental
\end{itemize}

\vspace{0.3cm}
\textbf{Soluciones:}
\begin{itemize}
\item Modelos de medidas repetidas
\item Modelos mixtos o jerárquicos
\item Rediseñar el experimento
\end{itemize}
\end{frame}

\section{Resumen}

\begin{frame}{Resumen de la Clase}
\textbf{ANOVA de dos factores:}
\begin{itemize}
\item Permite estudiar efectos de dos factores simultáneamente
\item Detecta interacciones entre factores
\item Tres pruebas: Factor A, Factor B, Interacción AB
\end{itemize}

\vspace{0.3cm}
\textbf{Diseños experimentales:}
\begin{itemize}
\item \textbf{DCA:} Simple, flexible, requiere homogeneidad
\item \textbf{DBCA:} Controla una fuente de variación, muy usado
\item \textbf{DCL:} Controla dos fuentes, muy restrictivo
\end{itemize}

\vspace{0.3cm}
\textbf{Verificación de supuestos:}
\begin{itemize}
\item Normalidad, homocedasticidad, independencia
\item Métodos gráficos y pruebas formales
\item Transformaciones y alternativas no paramétricas
\end{itemize}
\end{frame}

\begin{frame}{Próxima Clase}
\textbf{Clase 23: Implementación Computacional}

\vspace{0.3cm}
Veremos:
\begin{itemize}
\item ANOVA de dos factores en R y Python
\item Implementación de DCA, DBCA y DCL
\item Verificación de supuestos
\item Gráficos de interacción
\item Comparaciones múltiples post-hoc
\item Casos de estudio completos
\end{itemize}
\end{frame}

\end{document}
